\documentclass[11pt,oneside]{book}
\usepackage[paperwidth=7in,paperheight=9in,margin=0.72in]{geometry}
\usepackage{booktabs,tabularx,enumitem,xcolor,hyperref}
\definecolor{forest}{HTML}{356859}
\definecolor{gold}{HTML}{B8860B}
\hypersetup{colorlinks=true,linkcolor=forest,urlcolor=forest}
\setlength{\parindent}{0pt}
\setlength{\parskip}{0.45em}
\newcommand{\logday}[2]{\section{#1}\textbf{Aim.} #2\par\textbf{Result.}\hspace{0.5em}\rule{0.7\linewidth}{0.4pt}\par}
\title{Research Diary 2023\\\large A modular record for experiments and plans}
\author{Your Name\\Laboratory or project}
\date{Annual working copy}
\begin{document}
\frontmatter
\maketitle
\tableofcontents
\chapter*{Working rules}
Write the entry while the work is fresh. Keep a clear boundary between a
measurement and its interpretation, record deviations from the planned method,
and finish each session with a concrete next step. The compact page size is
intended for a long-running diary that can later be searched as one PDF.
\begin{enumerate}[leftmargin=*,label=\arabic*.]
\item State the aim and the relevant setup.
\item Record raw values before calculating a summary.
\item Link files, samples or instrument runs by a stable name.
\item Mark decisions and unresolved questions for the next session.
\end{enumerate}

\mainmatter
\chapter{January}
\logday{3 January --- instrument check}{Confirm the temperature controller and reference sample are ready for the week.}
\begin{tabularx}{\textwidth}{@{}lXrr@{}}
\toprule
Check & Condition & Reading & Status\\
\midrule
Warm-up & 20 minutes & 21.1\,\textdegree C & pass\\
Reference & Control sample & 1.018 & pass\\
Noise & 10-minute window & 0.012 & monitor\\
\bottomrule
\end{tabularx}
\logday{6 January --- comparison run}{Compare two preprocessing settings while holding the sample and operator constant.}
The smoothed signal is easier to interpret, but the raw trace remains the
authoritative record. Save both outputs with the same run identifier.
\logday{13 January --- project review}{Turn the first fortnight of observations into a decision for the next experiment.}
\begin{tabularx}{\textwidth}{@{}lXX@{}}
\toprule
Evidence & Decision & Follow-up\\
\midrule
Reference stable & Keep calibration procedure & Repeat at session end\\
Preprocessing changes peak height & Compare on held-out sample & Add blind replicate\\
\bottomrule
\end{tabularx}
\chapter{February}
\logday{2 February --- new series}{Start the next sample series and preserve the link to the January calibration.}
Use this chapter as the next month in the year. For a full annual diary,
duplicate the chapter and replace each example section with a dated entry.
\appendix
\chapter{Annual plan}
\begin{description}[leftmargin=2.4cm,style=sameline]
\item[Measure] Complete the reference check before every comparison.
\item[Review] Reserve one short entry each week for decisions and open questions.
\item[Archive] Keep source data and processed results under the same run identifier.
\end{description}
\end{document}
